\documentclass[conference]{IEEEtran}
\IEEEoverridecommandlockouts

\usepackage{cite}
\usepackage{amsmath,amssymb,amsfonts}
\usepackage{algorithmic}
\usepackage{graphicx}
\usepackage{textcomp}
\usepackage{xcolor}
\usepackage{booktabs}
\usepackage{url}

\begin{document}

\title{From Segmented Meshes to Queryable Semantic Digital Twins:\\
VLM-Driven Ontology Grounding for 3D Indoor Scenes}

\author{\IEEEauthorblockN{Anonymous Authors}
\IEEEauthorblockA{\textit{Institute} \\
City, Country \\
anonymous@example.org}
}

\maketitle

\begin{abstract}
TODO.
\end{abstract}

\begin{IEEEkeywords}
semantic digital twin, vision-language model, 3D scene understanding, ontology, knowledge graph
\end{IEEEkeywords}

\section{Introduction}
TODO.

\section{Related Work}
TODO.

\section{Overview}
TODO.

\section{Mesh Reconstruction and Segmentation}
\label{sec:recon}
TODO. (This stage is covered in detail in the companion paper~\cite{semdt_stage1}.)

\section{From Segmented Mesh to Semantic Digital Twin}
\label{sec:pipeline}

Given a segmented 3D mesh produced by the upstream reconstruction stage
(Section~\ref{sec:recon}), the goal of this stage is to construct a
\emph{semantic digital twin} (SemDT): a persistent, queryable
representation in which every segmented object is instantiated as a typed
entity in an ontology of indoor objects, with its geometry, pose, and
inter-object relations captured in a relational knowledge base.

The input is a collection of per-object triangle meshes together with an
inherited, partial taxonomy of \texttt{SemanticAnnotation} classes drawn
from the CRAM/\texttt{semantic\_digital\_twin} ontology. The output is an
instantiated \texttt{World} object, serialised to PostgreSQL via the
ORMatic object--relational mapper, in which each object is assigned a
concrete ontology class and the constructor-level dependencies between
annotations (e.g., a \texttt{Drawer} belongs to a \texttt{Cabinet}) are
resolved.

The stage is organised as a three-phase pipeline.

\subsection{Phase 1: Multi-View Class Extraction}
\label{sec:extract}

Classifying 3D objects directly from meshes is difficult for
general-purpose vision--language models (VLMs), which are trained
predominantly on 2D imagery. We therefore reduce the problem to a 2D
visual question answering task. The scene is rendered from three fixed
camera viewpoints (top, front-right, and front-left for Warsaw-style
scenes; auto-computed from the scene bounding box for HM3D). Objects are
processed in groups of $N$ (default $N{=}8$); within each group, the
target bodies are highlighted by overriding their mesh vertex colors
with a set of maximally distinct CSS3 colors, while the remaining scene
geometry is rendered in its original appearance to preserve spatial
context.

The three rendered views of the group are encoded as base64 PNG payloads
and submitted to a VLM (Qwen3-VL-30B-A3B-Instruct via OpenRouter) together
with: (i) a JSON schema listing the highlight color assigned to each
object, (ii) the closest CSS3 color name for disambiguation, and (iii)
the flattened list of available ontology labels from the taxonomy
export. The VLM is asked to return, for each highlighted object, a
predicted class label, a confidence score, and a short free-text
justification. Grouping objects reduces the number of VLM calls by an
order of magnitude relative to per-object querying while still allowing
the model to reason about each object in the context of its neighbours.

\subsection{Phase 2: Taxonomy-Aligned Label Refinement}
\label{sec:refine}

Raw VLM labels are lexically noisy: they may be synonyms of taxonomy
entries (\emph{fridge} vs.\ \emph{refrigerator}), morphological
variants (\emph{washer dryer} vs.\ \emph{washer-dryer}), or novel labels
that need to be introduced into the taxonomy. We resolve this in a
refinement step that uses sentence-level word embeddings to compute the
cosine similarity between each predicted label and every available
taxonomy class. Matches above an acceptance threshold are folded onto
the canonical class (e.g., \emph{fridge}~$\to$~\emph{refrigerator},
similarity 0.96); below-threshold predictions are preserved as candidate
new taxonomy entries for downstream review.

\subsection{Phase 3: Dependency Resolution and Persistence}
\label{sec:persist}

Many classes in the target ontology are not leaf concepts but structured
annotations whose constructors reference other annotations. For example,
instantiating a \texttt{Handle} requires a parent
\texttt{PrismaticConnection}; instantiating a \texttt{Drawer} requires a
\texttt{Cabinet} container. A naive instantiation order fails whenever a
required neighbour has not yet been created.

We address this with a three-step procedure. First, for each segmented
body we use the refined label from Phase~2 to look up the most specific
compatible \texttt{SemanticAnnotation} subclass in the taxonomy. Second,
we traverse the constructor signature of that class and, for each
unresolved argument, recursively instantiate or match an existing
annotation of the required type, calling a lightweight LLM
(Llama-3.3-70B-Instruct) only where the dependency cannot be satisfied
by geometric heuristics. Third, the resulting annotation graph is
topologically sorted and instantiated in dependency order, yielding a
valid \texttt{World} object. The world is persisted to PostgreSQL via
ORMatic so that it can be reloaded, rendered, and queried by downstream
reasoning components without re-running the VLM.

\subsection{Batch Execution}

The three phases are driven by a batch runner that iterates over all
rooms of a scene, caches VLM responses for reproducibility, records the
per-run taxonomy snapshot, and logs git commit hashes and model
identifiers into an experiment metadata file. This makes runs
reproducible across taxonomy evolutions: the same segmented mesh
re-processed against a later taxonomy yields a controlled set of
label changes attributable to the taxonomy rather than to VLM
non-determinism.

\section{Evaluation}
\label{sec:eval}

We evaluate the pipeline on scenes from the Habitat-Matterport 3D
(HM3D) semantic annotations dataset~\cite{hm3d} and, in future work,
on two real-world scenes (an instrumented kitchen and an office)
captured and reconstructed by the upstream stage.

\subsection{HM3D: Stage-2 Evaluation against Ground Truth}
\label{sec:eval_hm3d}

HM3D scenes provide per-face ground-truth (GT) labels in a
\texttt{.semantic.txt} lookup table, which allows us to evaluate the
classification component of the pipeline in isolation from the
reconstruction stage. We use three scenes from the HM3D \emph{minival}
split that cover a mix of residential floor plans and room counts:
\texttt{00802}, \texttt{00803}, and \texttt{00808}.

\paragraph{Protocol.}
For every segmented object, the pipeline produces a predicted label
which is aligned to the HM3D GT label vocabulary via the
embedding-based refinement of Section~\ref{sec:refine}. We report:
\emph{precision} and \emph{recall} of predicted labels against GT
labels; their harmonic mean (\emph{F1}); and a \emph{mean IoU}
(\emph{mIoU}) computed per class as the intersection-over-union of the
sets of objects assigned to that class in GT vs.\ prediction,
macro-averaged across classes present in either set. True positives,
false positives, and false negatives are counted at the
(room, object) level.

\paragraph{Results.}
Table~\ref{tab:hm3d} summarises the per-scene metrics. Across the three
scenes the pipeline achieves a mean precision of $0.834$, recall of
$0.676$, F1 of $0.746$, and mIoU of $0.534$. Precision is consistently
high ($>\!0.82$ on all scenes), indicating that when the VLM commits to
a label and the refinement step accepts it, the label is usually
correct. Recall is lower, reflecting (i) objects for which no taxonomy
class is a close enough match and the prediction is rejected by the
refinement threshold, and (ii) systematic confusions on visually
ambiguous or small classes (see error analysis below).

\begin{table}[t]
\centering
\caption{Stage-2 classification performance on three HM3D scenes.
TP / FP / FN are counted per (room, object); mIoU is the macro-averaged
per-class intersection-over-union of label assignments.}
\label{tab:hm3d}
\begin{tabular}{lccccc}
\toprule
Scene & Precision & Recall & F1 & mIoU & TP / FP / FN \\
\midrule
00802 & 0.847 & 0.675 & 0.751 & 0.575 & 714 / 129 / 344 \\
00803 & 0.826 & 0.635 & 0.718 & 0.459 & 530 / 112 / 305 \\
00808 & 0.829 & 0.718 & 0.770 & 0.567 & 597 / 123 / 234 \\
\midrule
\textbf{Mean} & \textbf{0.834} & \textbf{0.676} & \textbf{0.746} & \textbf{0.534} & --- \\
\bottomrule
\end{tabular}
\end{table}

\paragraph{Error analysis.}
Inspecting the per-class IoU breakdown on scene~\texttt{00808} reveals
three dominant failure modes. (i)~\emph{Taxonomy-level synonymy not
captured by embeddings}: labels such as \emph{poster} are matched
against \emph{picture} with similarity~$0.52$, below our acceptance
threshold, so they become false negatives on GT classes for which the
taxonomy does in fact carry a suitable class. (ii)~\emph{Visual
confusions on thin or occluded geometry}: the model confuses
\emph{cloth} with \emph{towel} (sim.~$0.66$) and \emph{handle} with
\emph{broom} when only a fragment of the object is visible in all three
views. (iii)~\emph{Fine-grained container classes}: \emph{kitchen
cabinet} is correctly retrieved only for $3.8\%$ of GT instances, most
cabinets being labelled with the coarser \emph{cabinet} (IoU~$0.12$).
Conversely, common and visually distinctive categories
(\emph{microwave}, \emph{refrigerator}, \emph{toilet},
\emph{dishwasher}, \emph{curtain}, \emph{television}) achieve IoU $=
1.0$ on all three scenes.

\paragraph{Scene 00800 (pilot).}
An earlier pilot run on scene~\texttt{00800} (59 objects) achieved an
object-level accuracy of $91.5\%$ ($54/59$ correct) using an earlier
taxonomy snapshot. These numbers are not directly comparable to
Table~\ref{tab:hm3d} because that run used per-object rather than
per-(room,object) counting and predated the embedding-based refinement,
and we report them here only for historical context.

\subsection{Ablations}
\label{sec:ablations}
TODO: report effect of number of camera viewpoints, group size $N$,
embedding acceptance threshold, and with/without prior-label conditioning.

\subsection{Real-World Scenes: End-to-End Evaluation}
\label{sec:eval_real}
TODO: full-pipeline evaluation on the IAI Kitchen and WTU Office
scenes, including reconstruction quality (Stage~1) and classification
quality (Stage~2) against a manually curated ground truth.

\subsection{Qualitative Results}
\label{sec:qualitative}
TODO: side-by-side renders of GT-coloured vs.\ predicted-coloured
scenes; taxonomy subtree visualisations; example queries against the
instantiated knowledge base.

\section{Discussion}
TODO.

\section{Conclusion}
TODO.

\bibliographystyle{IEEEtran}
\bibliography{references}

\end{document}
